\subsection*{本地保留七天，下游只保护一天}

设 E1 的本地恢复期限为七天，平台的逐键去重仅保留一天。取一条独立外部历史，平台余额初始 10000。第 0 天，A 保存 C1 与 \code{pay-C1}，平台完成 X1 并扣 3000，成功响应丢失。本地 C1 未决。第 1 天以后，平台删除 \code{pay-C1} 的去重记录。第 2 天，本地恢复者按原键再发，平台把它当作新支付，可能形成 X2，再扣 3000。

本地 S 和 C1 始终没有改变，外部效果却增加到两次。固定键只有在平台以同一身份解释所有可能生效的到达时，才提供相应上限。每次发送是否安全，要按平台当时对身份的解释规则判断。

\CTwelveFigure{ch12-original13}{本地映射长期存在，平台逐键窗口先结束，同一个 pay-C1 也可能再次首次执行}{fig:ch12-provider-ttl}

仅在平台窗口结束以前停止新发送，也要处理已经在途的副本。如果平台允许一个旧 HTTP 请求在排队很久以后才首次执行，而它已经忘记更早的成功，同样可能产生 X2。安全发送期限要结合传输和处理上界，或由平台在自己的受理点检查不可延长的截止，并为已受理工作保留身份保护直到终结。把最后一次发送在一天以内当成充分条件，会漏掉窗口以后的到达。

恢复任务因此保存下游能力和期限的适用依据。发送者每次读取原 C1，核对自己当前模式与平台是否保证这次可能生效的到达受保护。到达上界已经无法确认时，停止自动发送，T1 转入 \code{check}。查询和核对使用原平台范围及原支付身份，原输入保持。

\subsection*{一条可以终结未知的具体续路}

考虑支持原子关闭支付身份的平台。它为 \code{pay-C1} 提供 \code{SealPayment}：在处理原支付的同一仲裁域里，关闭以后新的首次受理资格，并返回原身份的完整结论或持久操作查询编号。这个能力是本例新增的平台前提，需要平台将关闭、支付受理与原结果查询接入同一身份仲裁。

若平台已经有成功 X1，关闭返回原成功证据，服务核对账户、金额、币种和原身份，执行 R。若平台已经受理但尚未完成，关闭返回 \code{accepted\_pending} 与原操作查询编号 H1，并保存禁止新首次受理的状态。C1/T1 以后只查询 H1，由平台持续保存这项已受理工作的结果，直到取得并保存终态。关闭请求自身超时，就按原 \code{pay-C1} 重复关闭或查询关闭状态，服务继续保持未决。

若平台返回 \code{closed\_without\_success}，这个结论要覆盖三个事实：过去没有成功，已受理或在途请求都不能再成功，未来也不能再首次受理。平台在同域封闭准入以后，把未开始的候选终结，并终结或等待已经受理的执行，最后才发布这份拒绝证据。

\CTwelveFigure{ch12-original14}{原子关闭把旧身份的过去与未来接到同一仲裁，未终结者沿持久编号继续查询}{fig:ch12-seal}

以平台一天的逐键窗口为例，恢复者在窗口结束前发起关闭，随后取得可长期查询的 H1 或终态证据。C1 保存关闭证据、原身份和 H1，T1 持续安排查询。若关闭服务不可达，恢复者在无法保证安全到达以前停止新支付发送，保留 T1 进入核查。平台保证的长期查询编号与关闭依据，给出了窗口结束以后的具体可靠续路。

平台若只有一天的逐键去重和普通查询，服务不能自行制造 \code{SealPayment}。此时可以停止发送，持续查询可解释的原支付账本，核对账单或进入有责任归属的人工流程。查到并核验 X1 可以确认成功；查空继续表示未知。要确认最终未付款，就取得完整账本覆盖、全部在途受理已经结束以及禁止未来成功的依据。证据一时不可得，C1 继续未决，系统通过限制新增受理量管理积压。

\CTwelveFigure{ch12-final-evidence}{最终拒绝的证据覆盖过去、已受理与在途、未来三段，普通查空只有当前读取结果}{fig:ch12-final-evidence}

对于新购买，准入服务提前判断平台能否覆盖预期的发送与恢复路径。无法提供所要求保证的平台，可以只接受另一种明确公布的处理方式，或暂停相关支付方式。已经受理的 C1 留在核查链中，停止新受理帮助控制新增责任。

\subsection*{过期响应准确说出入口发生了什么}

E1 到期以后，购买入口可以返回下面的响应。\code{operation\_expired} 描述身份资格，\code{outcome} 则表达该入口没有据此得出资金结论。示例查询位置在原 C1 已有且调用者经过授权时给出：

\begin{BookCode}
{
  "code": "operation_expired",
  "session_id": "E1",
  "idempotency_key": "K1",
  "retry_until": "2026-10-18T10:00:00Z",
  "outcome": "unknown_at_this_entry",
  "purchase_id": "P1",
  "status_url": "/v1/purchases/P1"
}
\end{BookCode}

\code{unknown\_at\_this\_entry} 的含义是这个过期响应没有声明购买的成功或失败，已知 P1 的业务查询可以给出更多证据。如果会话记录及逐笔定位都已经回收，服务无法确定 P1，就省去 \code{purchase\_id} 与具体状态地址，提供经账户授权的查询入口。缺失信息用相应响应结构表达。

业务查询可以返回 P1 的已知成功和原结果，也可以返回当前核查进度。它使用自己的授权、保留范围与读取一致性规则。查询入口按 P1 或原支付证据定位，旧购买提交入口按 E1 决定能否恢复或首次受理，二者在协议中各有明确用途。

\CTwelveFigure{ch12-original15}{过期只结束旧入口，成功证据结束原购买，未知状态进入核查，新意图取得新身份}{fig:ch12-expired-options}

\subsection*{新操作需要新的业务决定}

Alice 查到原 P1 已成功，却确实还想再买一份，可以明确确认另一项购买。服务登记新 E2，客户端持久保存新 K2 与当次报价，形成新的完整 S2。S2 与原 S1 可以各完成一次，这是两项业务。新报价、授权与业务限制在 S2 首次受理时判断，原 P1 的结果继续按原事实解释。

另一种情况是原 C1 未决，Alice 希望再试一次以保证买到。客户端如果直接把旧任务复制成 E2/K2，平台便可能分别完成旧 C1 与新 C2。需要一份的业务意图因重建身份而得到两份购买。界面应把原未决购买和新购买确认分开，显示原关联与核查状态。只有业务明确允许承担两项各自可能成功的购买，才提交第二项。

如果业务目标是账户对某一预约只可成功一次，可以为该预约另设长期唯一的业务事实身份，并在新会话受理时仲裁它。这样 E1/K1 和 E2/K2 可以是不同请求身份，却引用同一预约事实。其长期唯一索引及输入规则承担了跨会话合并责任，存储成本也跟着业务生命周期保留。对任意可重复购买的商品，单用商品编号作唯一键会把 Alice 的两次合法购买错误合并。

最后把回收过程接回本章开头的快照。活动 E1 的快照保存未确认的完整结果、已确认的 $M_{\min}$/U、确认水位和固定资格。关闭后的快照可以省去已安全回收的逐笔结果与适配，保留不会重用旧代次的登记进度，并保留尚未终结的 C/T 及其证据关系。恢复者对缺失 E1 拒绝旧购买，对已受理 C1 继续核查。每一次删除，都由另一份足以决定后续动作的信息接手。

逐操作身份的成本由此缩到了明确范围：完整结果保留到确认或期限，活动会话内的任意键依据保留到关闭，未完成责任保留到有证据的终结。有些业务更新能把重复输入吸收进当前状态，另一些更新每次到达都会产生新效果。要再减少逐笔历史，就要检查业务操作本身允许怎样重复、交错与合并。
\FloatBarrier
